\section{Conclusiones}

Se desarrolló una PINN inversa para estimar estados neurovasculares, una HRF operacional y parámetros efectivos del modelo Balloon--Windkessel. La reformulación que separa el entrenamiento de la predicción mediante integración ODE independiente fue fundamental para evitar soluciones físicamente blandas con baja generalización.

En 60 experimentos sintéticos, la PINN recuperó la señal, la HRF y los parámetros con alta precisión incluso bajo ruido elevado, y superó significativamente al GLM y a la MLP. Estos resultados validan la metodología en un entorno con ground truth conocido.

En los datos reales del sujeto HCP 100206, la PINN obtuvo el menor RMSE medio y ganó cuatro de ocho evaluaciones. Las diferencias de error no fueron estadísticamente significativas, pero PINN y MLP presentaron una correlación temporal superior al GLM. La MLP mostró la mayor brecha de generalización, mientras que la PINN produjo HRF altamente consistentes en forma.

En conjunto, la incorporación de conocimiento fisiológico mejoró la identificación y generalización en datos controlados y ofreció una representación interpretable en fMRI real. Sin embargo, la estimación de parámetros individuales, en especial $\tau$, requiere más datos y análisis de incertidumbre antes de atribuir significado fisiológico o clínico directo.
